% ============================================================================
% Supplemental Material
% "Phonon scattering by a disclination loop: forward decoherence
%  and a flat Boltzmann transport integrand"
% ============================================================================
\documentclass[11pt]{article}
\usepackage{amsmath,amssymb}
\usepackage{graphicx}
\usepackage{hyperref}
\usepackage[margin=1in]{geometry}

\begin{document}

\begin{center}
{\large\bfseries Supplemental Material}\\[6pt]
{\itshape Phonon scattering by a disclination loop: forward decoherence
and a flat Boltzmann transport integrand}\\[4pt]
Alexandru Toader
\end{center}
\medskip

% ====================================================================
\section{FDTD simulation details}
\label{sec:fdtd}

The scattering problem is solved by finite-difference time-domain (FDTD)
integration of the lattice equation of motion
%
\begin{equation}
m\ddot{u}_i = \sum_j K_{ij}(u_j - u_i)
\end{equation}
%
on a simple cubic lattice with $K_1 = 1.0$, $K_2 = 0.5$ ($K_1 = 2K_2$ for
isotropic sound speed $c = \sqrt{K_1 + 4K_2} = \sqrt{3}$). The time step is
$\Delta t = 0.25$, well within the CFL stability limit.

\textbf{Peierls gauge.}
Bonds crossing the Dirac disk of the disclination loop are modified:
$u_j \to \mathcal{R}(2\pi\alpha)\,u_j$, where $\mathcal{R}$ is a rotation
matrix of angle $2\pi\alpha$ in the $(x,y)$ plane. This produces a coupling
$\delta K = K_1(\mathcal{R} - I)$ on each gauged bond. Since $\mathcal{R}$
acts only on $(u_x, u_y)$, the $u_z$ component decouples:
$N_{\rm pol} = 2$ independent polarizations contribute to transport.

\textbf{Boundary conditions.}
Perfectly matched layers (PML) of width $D_W = 15$ lattice spacings with
damping coefficient $\sigma_s = 1.5$ absorb outgoing waves on all faces. PML
convergence is verified: the spread across $D_W = 10, 15, 20$ is less than
2\% at all $k \geq 0.5$. The $k = 0.3$ mode is sensitive to box size rather
than PML (see below).

\textbf{Wave packet.}
A Gaussian wave packet with spatial width $\sigma_x = 8$ is launched along
$+x$ at wavenumber~$k$. Bandwidth independence is verified: $\sigma_x = 6,
8, 12$ give results within 5\% at $k = 0.9$ and $k = 1.5$.

\textbf{Box size.}
The default box is $L = 100$ lattice spacings. Convergence is verified by
comparison with $L = 120$: agreement within 5\% at all~$k$. At $L = 80$,
the measurement sphere is too close to the PML and $k = 0.3$ overestimates
$\sigma_{\rm tr}$ by 44\%. At $k \geq 0.9$, $L = 80$ bias is below 15\%.

\textbf{Measurement.}
Scattered waves are recorded on a sphere of radius $r_m = 20$ surrounding the
defect, centered at the loop center. The recording starts at
$x_{\rm start} = L/2 - r_m$ to capture the wave packet as it enters the
measurement sphere. The transport cross-section is
%
\begin{equation}
\sigma_{\rm tr} = \frac{1}{P_{\rm inc}}
  \int (1 - \cos\theta_s)\,\frac{d\sigma}{d\Omega}\,d\Omega
\end{equation}
%
where $\theta_s$ is the scattering angle relative to the incident direction
$+x$, and $P_{\rm inc}$ is the incident power.

% ====================================================================
\section{Transport formula derivation}
\label{sec:transport}

The Boltzmann phonon drag coefficient for a single defect of radius~$R$ in a
Debye solid is
%
\begin{equation}
\kappa = \frac{R}{2\pi^2} \int_0^{k_{\rm max}}
  \sin^2(k)\,\sigma_{\rm tr}(k)\,dk \,.
\end{equation}

\textbf{Derivation.}
Start from the standard drag force per unit thermal gradient:
%
\begin{equation}
F_{\rm drag} = \frac{1}{V} \sum_{\mathbf{k},s}
  \hbar\omega\,v_g\,\tau^{-1}\,\frac{\partial n}{\partial T}\,\sigma_{\rm tr}
\end{equation}
%
In the Debye approximation with lattice dispersion
$\omega(k) = 2c\sin(k/2)$, the density of states along one cubic axis is
%
\begin{equation}
g(k) = \frac{k^2}{2\pi^2} \to \frac{\sin^2(k)}{2\pi^2}
\end{equation}
%
after the change of variable from continuum~$k$ to lattice dispersion (the
Jacobian $dk/d\omega = 1/v_g$ combined with the Boltzmann weight~$v_g$
produces $\sin^2(k)$ from $\omega^2/c^2$). The resulting integral has no tail
extrapolation: we sample 7 wavenumbers $k \in [0.3, 1.5]$ covering the
dominant part of the Brillouin zone.

\textbf{Polarization.}
Since $\mathcal{R}(2\pi\alpha)$ acts only on $(u_x, u_y)$, the $u_z$
polarization is unscattered. Two transverse polarizations contribute:
$N_{\rm pol} = 2$.

% ====================================================================
\section{Systematics}
\label{sec:systematics}

The dominant systematic uncertainties are the gauging prescription (NN vs NNN)
and the $k$-space cutoff.

\subsection{NN vs NNN gauging}

The Peierls gauge can be applied to nearest-neighbor bonds only (NN) or also
to next-nearest-neighbor bonds (NNN). Both gaugings produce identical holonomy
$\mathcal{R}(2\pi\alpha)$ around the disclination, but distribute the gauge
differently at the lattice scale.

NNN gauging gives systematically higher $\sigma_{\rm tr}$ at all~$k$ and
produces a less flat integrand (CV~$= 29\%$ at $\alpha = 0.30$ vs 5.6\% for
NN). The physical reason: NN $z$-bonds have $\Delta x = 0$, so the incident
phase $\exp(ikx)$ is identically~1 at all~$k$, making the per-bond vertex
$k$-independent. NNN bonds have $\Delta x = 1$, introducing $k$-dependent
phases.

At $\alpha = 0.30$, $k \leq 1.5$: $\kappa_{\rm NN} = 0.844$,
$\kappa_{\rm NNN} = 2.227$. The NNN gauging overshoots $\kappa = 1$.

\subsection{\texorpdfstring{$k$}{k}-cutoff}

The integral is evaluated at two cutoffs: $k \leq 0.9$ (conservative) and
$k \leq 1.5$ (full range). The ratio is $1.9$--$2.7\times$ depending on
$\alpha$.

\subsection{Gauge invariance}

Rotating the Dirac disk orientation (gauge choice) changes $\sigma_{\rm tr}$
by less than 2.6\% at all~$k$, confirming that the result is a property of
the disclination topology, not the gauge surface.

\subsection{Other systematics}

\begin{table}[h]
\begin{tabular}{ll}
\hline
Source & Effect \\
\hline
Angular grid ($N_\theta = 13$, $N_\phi = 24$) & Converged within 5\% \\
Near-field correction ($r_m = 20$) & Small ($< 3\%$) \\
$k$/$k_{\rm eff}$ phase mismatch & $< 9\%$ at $k = 1.5$ \\
$K_1/K_2$ variation & $\kappa = 0.75$--$1.23$ \\
\hline
\end{tabular}
\end{table}

The $T$-matrix calculation uses the lattice wavenumber~$k$ for the incident
phase and $k_{\rm eff} = 2\sin(k/2)$ for the outgoing phase (from the
continuum Green function). The mismatch $k - k_{\rm eff} = O(k^3/24)$
reaches $\sim 9\%$ at $k = 1.5$. This is validated empirically: the MS
$T$-matrix reproduces 75--110\% of the FDTD exponent correction at $R = 5$.

\subsection{Uncertainty budget}

\begin{table}[h]
\begin{tabular}{ll}
\hline
Systematic & Factor \\
\hline
NN vs NNN gauging & $2.1$--$4.9\times$ \\
$k$-cutoff & $1.9$--$2.7\times$ \\
Gauge invariance & $< 2.6\%$ \\
Angular grid & $< 5\%$ \\
PML width & $< 2\%$ \\
Box size ($L = 100$ vs $120$) & $< 5\%$ \\
Bandwidth ($\sigma_x$) & $< 5\%$ \\
$k$/$k_{\rm eff}$ phase mismatch & $< 9\%$ at $k = 1.5$ \\
$K_1/K_2$ & $\kappa = 0.75$--$1.23$ \\
\hline
\end{tabular}
\end{table}

% ====================================================================
\section{Per-bond vertex \texorpdfstring{$V(k)$}{V(k)}}
\label{sec:vertex}

The Born vertex
$V(k) = c_{m1}^2 Z_{\rm mono}(k) + s_\phi^2 Z_{\rm dipo}(k)$ combines a
monopole angular integral $Z_{\rm mono} = 8\pi(1 + \sin k/k)$ and a dipole
integral $Z_{\rm dipo} = 8\pi(1 - \sin k/k)$. The monopole term is
approximately constant ($Z_{\rm mono}$ CV~$= 5.9\%$), while the dipole varies
strongly ($Z_{\rm dipo}$ CV~$= 72.7\%$). At strong coupling
($\alpha \geq 0.25$), $|c_{m1}|^2 \gg s_\phi^2$ and the monopole dominates,
giving $V(k) \approx \mathrm{const}$.

$V(k)$ at $\alpha = 0.30$:

\begin{table}[h]
\begin{tabular}{lccccccc}
\hline
$k$    & 0.3   & 0.5   & 0.7   & 0.9   & 1.1   & 1.3   & 1.5 \\
\hline
$V(k)$ & 85.83 & 85.30 & 84.51 & 83.50 & 82.27 & 80.87 & 79.32 \\
\hline
\end{tabular}
\end{table}

CV~$= 2.7\%$. The variation is less than 8\% peak-to-peak across the sampled
range. At $\alpha = 0.25$, $|c_{m1}| = |s_\phi|$ exactly, making
$V(k) = (c_{m1}^2 + s_\phi^2)(Z_{\rm mono} + Z_{\rm dipo})/2 =
\mathrm{const}$ (CV~$= 0.0\%$, since
$Z_{\rm mono} + Z_{\rm dipo} = 16\pi$).

$V(k)$ CV vs $\alpha$:

\begin{table}[h]
\begin{tabular}{lcccccc}
\hline
$\alpha$ & 0.05   & 0.10   & 0.20  & 0.25  & 0.30  & 0.50 \\
\hline
CV       & 54.1\% & 28.2\% & 4.6\% & 0.0\% & 2.7\% & 5.9\% \\
\hline
\end{tabular}
\end{table}

The vertex constancy has a sharp threshold near $\alpha = 0.25$ (the
monopole-dipole crossover), with CV~$< 6\%$ for $\alpha \geq 0.20$.

% ====================================================================
\section{Assembly chain details}
\label{sec:assembly}

The flat integrand factorizes as
%
\begin{multline}
\sin^2(k) \cdot \sigma_{\rm ring} =
  \underbrace{4\sin^2(k/2)\cos^2(k/2)}_{\sin^2(k)} \\
  \cdot \underbrace{\frac{C_0\,V(k)}{\cos^2(k/2)}}_{1/v_g^2}
  \cdot N_{\rm eff}(k)
\end{multline}

\textbf{Step~1: Algebraic cancellation.}
$\cos^2(k/2)$ from $\sin^2(k)$ cancels $\cos^2(k/2)$ from $1/v_g^2$,
leaving $4\sin^2(k/2) \cdot V(k) \cdot N_{\rm eff}(k)$.

\textbf{Step~2: Constant vertex.}
At $\alpha \geq 0.25$, the monopole coupling $|c_{m1}|^2 Z_{\rm mono}$
dominates the vertex. Since $Z_{\rm mono}$ is approximately constant
(CV~$< 7\%$), $V(k) \approx \mathrm{const}$ (CV~$< 6\%$).

\textbf{Step~3: $N_{\rm eff}$ balance.}
The remaining factor $4\sin^2(k/2) \cdot N_{\rm eff}$ must be approximately
constant, requiring $N_{\rm eff} \propto 1/\sin^2(k/2)$, i.e., exponent
$\approx -2.0$.

The Born exponent is $-5/2$: the filled disk forward cone has angular widths
$\delta \sim 1/\sqrt{kR}$, $\eta \sim 1/(kR)$, giving
$\sigma_{\rm tot} \sim (kR)^{-3/2}$; the transport weight
$(1 - \cos\theta)$ adds $-1$. Multiple scattering shifts this to
$\approx -2.0$ (shift $+0.45$ at $\alpha = 0.30$).

The lattice form $1/\sin^2(k/2)$ fits better than the continuum $1/k^2$:
CV~$= 5.8\%$ vs 7.1\%. This lattice dispersion correction reflects the
$\sin(k/2)$ appearing in the exact dispersion relation.

\textbf{Balance point.}
The integrand is flattest at $\alpha \approx 0.25$ (CV~$= 3.7\%$). At
$\alpha = 0.30$: CV~$= 5.6\%$. The window $\alpha \in [0.20, 0.40]$ gives
CV~$< 11\%$.

% ====================================================================
\section{FDTD convergence}
\label{sec:convergence}

Box size convergence at $R = 5$, $\alpha = 0.30$. All three box sizes are
run with identical FDTD parameters; $L = 120$ serves as reference.

\begin{table}[h]
\begin{tabular}{ll}
\hline
Test & Result \\
\hline
$L = 100$ vs $L = 120$ & Agreement within 5\% at all $k$ \\
$L = 80$, $k = 0.3$ & Overestimates by 44\% (PML too close) \\
$L = 80$, $k \geq 0.9$ & Bias $< 15\%$ \\
Integrand CV at $L = 100$ & 5.6\% (converged) \\
\hline
\end{tabular}
\end{table}

Reference $\sigma_{\rm tr}(k)$ values at $L = 100$, $R = 5$,
$\alpha = 0.30$:

\begin{table}[h]
\begin{tabular}{lccccccc}
\hline
$k$ & 0.3 & 0.5 & 0.7 & 0.9 & 1.1 & 1.3 & 1.5 \\
\hline
$\sigma_{\rm tr}$ & 28.3 & 12.3 & 7.09 & 4.69 & 3.50 & 2.91 & 2.60 \\
$\sin^2(k)\,\sigma_{\rm tr}$ & 2.47 & 2.83 & 2.94 & 2.88 & 2.78 & 2.70 & 2.59 \\
\hline
\end{tabular}
\end{table}

% ====================================================================
\section{Negative tests}
\label{sec:negative}

The flat integrand is not a generic property. It fails under the following
modifications:

\textbf{Born approximation only.}
Without multiple scattering, the Born integrand has CV~$> 30\%$
(approximately $19\times$ variation across the BZ). The Born exponent $-5/2$
makes low-$k$ modes dominate.

\textbf{Weak coupling.}
At $\alpha = 0.05$, the per-bond vertex $V(k)$ varies strongly
(CV~$= 54\%$) because the dipole term $|s_\phi|^2 Z_{\rm dipo}$ dominates
over the monopole. At $\alpha = 0.10$, the integrand CV exceeds 15\%.

\textbf{NNN gauging.}
With NNN Peierls gauge, the integrand CV~$= 29\%$ at $\alpha = 0.30$. The
$k$-dependent phases $\exp(ik\Delta x)$ from NNN bonds break the per-bond
vertex constancy.

\textbf{Aharonov--Bohm formula.}
The AB cross-section $\sigma_{\rm AB} \sim 1/k$ gives an integrand
CV~$> 20\%$, amplitude $6.5\times$ below FDTD, and does not reproduce the
observed $k$-dependence, $\alpha$-dependence, or $R$-scaling.

These negative tests establish that the flat integrand requires: (i)~multiple
scattering ($N_{\rm eff}$ exponent shift), (ii)~strong coupling
($V(k) \approx \mathrm{const}$), and (iii)~NN gauging (per-bond phase
independence).

% ====================================================================
\section{\texorpdfstring{$\kappa(\alpha)$}{kappa(alpha)} tables}
\label{sec:kappa}

\subsection{NN gauging}

\begin{table}[h]
\begin{tabular}{lcc}
\hline
$\alpha$ & $\kappa$ ($k \leq 0.9$) & $\kappa$ ($k \leq 1.5$) \\
\hline
0.10 & 0.022 & 0.049 \\
0.15 & 0.071 & 0.145 \\
0.20 & 0.162 & 0.315 \\
0.25 & 0.287 & 0.557 \\
0.30 & 0.428 & 0.844 \\
0.40 & 0.674 & 1.386 \\
0.50 & 0.769 & 1.614 \\
\hline
\end{tabular}
\end{table}

\subsection{NNN gauging}

\begin{table}[h]
\begin{tabular}{lcc}
\hline
$\alpha$ & $\kappa$ ($k \leq 0.9$) & $\kappa$ ($k \leq 1.5$) \\
\hline
0.10 & 0.087 & 0.238 \\
0.15 & 0.209 & 0.556 \\
0.20 & 0.393 & 1.016 \\
0.25 & 0.628 & 1.593 \\
0.30 & 0.889 & 2.227 \\
\hline
\end{tabular}
\end{table}

$\kappa = 1$ crossing: $\alpha \approx 0.33$ (NN, $k \leq 1.5$); with NNN
gauging the crossing shifts to $\alpha \approx 0.20$. The interval
$\alpha \approx 0.20$--$0.33$ is bracketed by gauging prescription, not
$k$-cutoff (at $k \leq 0.9$, NN never reaches $\kappa = 1$).

% ====================================================================
\section{Forward decoherence: microscopic structure}
\label{sec:forward}

The main text reports that multiple scattering shifts the Born exponent from
$-5/2$ to $\approx -2$ through forward decoherence. This section details
the microscopic structure of the mechanism.

\subsection{$T$-matrix decomposition}

The $T$-matrix $T = (I - VG)^{-1}V$ modifies the scattering amplitude in two
ways: (i)~diagonal terms $T_{jj}$ enhance the per-bond amplitude
($|T_{jj}/V|^2 = 1.35$), and (ii)~off-diagonal cross terms between bonds
$j \neq l$ produce destructive interference (cross power $= -137$ at
$k = 0.3$, which is 73\% of the diagonal contribution). The net effect is
suppression at low~$k$ (MS/Born power ratio $= 0.61$).

At strong coupling ($\alpha = 0.30$), the cross term is 73\% of the diagonal;
at weak coupling ($\alpha = 0.05$), it is only 2.8\%. The cross-term
strength scales as $26\times$ from $\alpha = 0.05$ to $\alpha = 0.30$. This
is why weak coupling shows no forward decoherence.

\subsection{Eigenvalue structure}

The eigenvalues of $VG$ control the resummation. At $k = 0.3$:
$|\lambda_{\rm max}|/|\lambda_2| = 2.52$---a single dominant mode concentrates
the cooperation, justifying a mean-field picture. At $k = 1.5$: the ratio drops
to $1.01$ (no dominant mode; uniform spectrum). This transition underlies the
$k$-dependent decoherence.

The eigenvalue magnitudes are bounded: $|\lambda_{\rm max}| = 0.3$--$0.86$ for
$R \leq 9$, confirming that the system is not resonant ($|\lambda| < 1$) and
the Born series converges.

\subsection{Path excess and cooperation range}

Each off-diagonal $T$-matrix contribution involves a propagation path
$i \to l \to j$ with path excess $\mathrm{pe} = r_{il} + r_{lj} - r_{ij}$. The
raw geometric distribution has median $\mathrm{pe} = 4.0$ lattice spacings (not
short-range). However, the Green function weighting $1/(r_{il}\,r_{lj})$
concentrates the effective distribution at short range: the $1/r^3$-weighted
median is $\mathrm{pe} = 2.0$, and 62.5\% of the weighted sum comes from
triplets with $\mathrm{pe} < 3.0$ (vs.\ 39.5\% by count).

This explains the convergence with propagation cutoff: physical (oscillating)
Green functions converge at $r_{\rm cut} = 3$, while static $1/r$ does not
converge even at $r_{\rm cut} = 5$. The cooperation is local---not because
the geometry favors short paths, but because the $1/r$ weighting in the Green
function suppresses long-range contributions.

\subsection{$R$-scaling of the shift}

The transport exponent shift $\Delta p = p_{\rm MS} - p_{\rm Born} \approx
+0.45$ is independent of disk radius: $\Delta p = +0.41$ ($R = 3$), $+0.45$
($R = 5$), $+0.40$ ($R = 7$), $+0.41$ ($R = 9$), with CV~$= 4.2\%$. The shift
is proportional to coupling strength: $\Delta p = C\,|V|$ with
$C = 0.334 \pm 0.012$ (CV~$= 3.7\%$ over $\alpha = 0.15$--$0.40$).

Despite the $R$-independent shift, the absolute $\sigma_{\rm tr}$ grows with
$R$. Born: $\sigma_{\rm tr}^{\rm Born} \sim R^{3/2}$ (from stationary phase).
Multiple scattering slightly reduces the $R$-exponent ($p_{\rm MS} < p_{\rm
Born}$) because forward decoherence strengthens with disk size. FDTD gives
$\sigma_{\rm tr} \sim R^{1.63}$, consistent with the Born prediction plus
sub-leading corrections.

\subsection{Crossover wavenumber}

The crossover $k^*$ where the MS/Born power ratio crosses unity grows with
disk radius: $k^* = 0.69$ ($R = 3$), $0.83$ ($R = 5$), $1.15$ ($R = 7$),
$1.20$ ($R = 9$), $1.56$ ($R = 12$). The scaling follows $k^* \sim R^{0.59}$.
At $R \geq 25$, the crossover reaches the BZ edge ($k^* > \pi$), meaning the
entire spectrum is suppressed---consistent with the Born series eventually
diverging at large~$R$.

\subsection{$N_{\rm eff}$ validation}

The MS calculation uses a scalar Green function approximation, producing
$\sigma_{\rm tr}$ values $\sim\!14$--$15\times$ larger than FDTD (weakly
$k$-dependent). This gap arises entirely from the per-bond normalization
(scalar vs.\ vector field): the
collective $N_{\rm eff} = \sigma_{\rm ring}/\sigma_{\rm bond}$ from MS matches
FDTD with mean ratio $1.05$ (CV~$= 14\%$). The mechanism correctly predicts
how $N$ bonds interact, even though the single-bond cross-section has an overall
normalization offset. This normalization difference does not affect the exponent
or $N_{\rm eff}$ scaling, which are the quantities relevant to the mechanism.

\subsection{Integrand flatness and coupling}

The MS integrand CV decreases monotonically with coupling strength: 39.5\%
($\alpha = 0.20$), 35.1\% ($\alpha = 0.25$), 31.5\% ($\alpha = 0.30$), 28.8\%
($\alpha = 0.35$), 26.9\% ($\alpha = 0.40$). There is no minimum at
$\alpha = 0.25$ despite the per-bond vertex $V(k)$ being exactly constant
there. Stronger coupling gives a larger shift and always flattens the integrand
further. (The full FDTD integrand is flattest at $\alpha \approx 0.25$
(\S\ref{sec:assembly}), where $V(k)$ is exactly constant; the MS integrand
reported here captures only the $N_{\rm eff}$ contribution and does not
include the per-bond vertex variation.)

\end{document}
